\documentclass[11pt,a4paper]{article}

\usepackage[margin=18mm]{geometry}
\usepackage{fontspec}
\usepackage{xeCJK}
\usepackage{graphicx}
\usepackage{subcaption}
\usepackage{booktabs}
\usepackage{array}
\usepackage{float}
\usepackage{xcolor}
\usepackage{hyperref}
\usepackage{caption}
\usepackage{fancyhdr}

\setmainfont{Times New Roman}
\setCJKmainfont{Microsoft JhengHei}
\hypersetup{colorlinks=true, linkcolor=black, urlcolor=blue}
\pagestyle{fancy}
\fancyhf{}
\lhead{Mindlin 板屈曲模態測試}
\rhead{\thepage}
\renewcommand{\headrulewidth}{0.4pt}
\captionsetup{font=small, labelfont=bf}
\setlength{\parindent}{0pt}
\setlength{\parskip}{0.55em}

\graphicspath{{../../vtk/}}

\title{\textbf{Mindlin 板屈曲模態與網格比較}\\[0.35em]
\large Q4 與 T3 元素之模態圖整理}
\author{}
\date{}

\begin{document}
\maketitle

\section{模型與測試摘要}

本報告整理 Mindlin 板屈曲算例中，四節點四邊形元素（Q4）與三節點三角形元素（T3）所得之模態圖。圖中 M1--M4 表示輸出的四個候選模態，色階為正規化後的橫向位移 $w$。

\begin{table}[H]
\centering
\caption{模型參數與靜力一致性測試}
\begin{tabular}{>{\raggedright\arraybackslash}p{0.28\textwidth}p{0.62\textwidth}}
\toprule
項目 & 數值或說明 \\
\midrule
材料與幾何參數 & $E=200\times 10^9$, $\nu=0.3$, $h=1.0$, $a=b=1.0$ \\
彎曲剛度 & $D^b = Eh^3/[12(1-\nu^2)]$ \\
理論屈曲係數 & $k_{\mathrm{exact}}=4.0$ \\
測試網格 & $16\times16$ \\
靜力測試載重參數 & $\lambda_{\mathrm{test}}=7.230479\times10^{10}$ \\
相對誤差 & $\|d-d_{\mathrm{exact}}\|/\|d_{\mathrm{exact}}\|=1.198424\times10^{-14}$ \\
$w$ 相對誤差 & $4.326165\times10^{-15}$ \\
$\phi$ 相對誤差 & $1.221666\times10^{-14}$ \\
\bottomrule
\end{tabular}
\end{table}

靜力一致性測試結果接近雙精度機器誤差，表示目前離散矩陣拼接、自由度排列以及 $w,\phi_1,\phi_2$ 的輸出映射在代數層面是自洽的。後續模態圖中的震盪現象，較可能與元素型式、模態篩選方式、幾何剛度矩陣之奇異性或鎖定效應有關。

\clearpage

\section{Q4 元素模態圖}

\begin{figure}[H]
\centering
\begin{subfigure}{0.48\textwidth}
  \centering
  \includegraphics[width=\linewidth]{m1x_q4.png}
  \caption{M1\_Q4}
\end{subfigure}\hfill
\begin{subfigure}{0.48\textwidth}
  \centering
  \includegraphics[width=\linewidth]{m2x_q4.png}
  \caption{M2\_Q4}
\end{subfigure}

\vspace{0.8em}

\begin{subfigure}{0.48\textwidth}
  \centering
  \includegraphics[width=\linewidth]{m3x_q4.png}
  \caption{M3\_Q4}
\end{subfigure}\hfill
\begin{subfigure}{0.48\textwidth}
  \centering
  \includegraphics[width=\linewidth]{m4x_q4.png}
  \caption{M4\_Q4}
\end{subfigure}
\caption{Q4 元素所得 M1--M4 模態圖。}
\end{figure}

\clearpage

\section{T3 元素模態圖}

\begin{figure}[H]
\centering
\begin{subfigure}{0.48\textwidth}
  \centering
  \includegraphics[width=\linewidth]{m1x_t3.png}
  \caption{M1\_T3}
\end{subfigure}\hfill
\begin{subfigure}{0.48\textwidth}
  \centering
  \includegraphics[width=\linewidth]{m2x_t3.png}
  \caption{M2\_T3}
\end{subfigure}

\vspace{0.8em}

\begin{subfigure}{0.48\textwidth}
  \centering
  \includegraphics[width=\linewidth]{m3x_t3.png}
  \caption{M3\_T3}
\end{subfigure}\hfill
\begin{subfigure}{0.48\textwidth}
  \centering
  \includegraphics[width=\linewidth]{m4x_t3.png}
  \caption{M4\_T3}
\end{subfigure}
\caption{T3 元素所得 M1--M4 模態圖。}
\end{figure}

\clearpage

\section{同一模態下 Q4 與 T3 對照}

\begin{figure}[H]
\centering
\begin{subfigure}{0.48\textwidth}
  \centering
  \includegraphics[width=\linewidth]{m1x_q4.png}
  \caption{M1\_Q4}
\end{subfigure}\hfill
\begin{subfigure}{0.48\textwidth}
  \centering
  \includegraphics[width=\linewidth]{m1x_t3.png}
  \caption{M1\_T3}
\end{subfigure}
\caption{M1 模態之 Q4 與 T3 對照。}
\end{figure}

\begin{figure}[H]
\centering
\begin{subfigure}{0.48\textwidth}
  \centering
  \includegraphics[width=\linewidth]{m2x_q4.png}
  \caption{M2\_Q4}
\end{subfigure}\hfill
\begin{subfigure}{0.48\textwidth}
  \centering
  \includegraphics[width=\linewidth]{m2x_t3.png}
  \caption{M2\_T3}
\end{subfigure}
\caption{M2 模態之 Q4 與 T3 對照。}
\end{figure}

\clearpage

\begin{figure}[H]
\centering
\begin{subfigure}{0.48\textwidth}
  \centering
  \includegraphics[width=\linewidth]{m3x_q4.png}
  \caption{M3\_Q4}
\end{subfigure}\hfill
\begin{subfigure}{0.48\textwidth}
  \centering
  \includegraphics[width=\linewidth]{m3x_t3.png}
  \caption{M3\_T3}
\end{subfigure}
\caption{M3 模態之 Q4 與 T3 對照。}
\end{figure}

\begin{figure}[H]
\centering
\begin{subfigure}{0.48\textwidth}
  \centering
  \includegraphics[width=\linewidth]{m4x_q4.png}
  \caption{M4\_Q4}
\end{subfigure}\hfill
\begin{subfigure}{0.48\textwidth}
  \centering
  \includegraphics[width=\linewidth]{m4x_t3.png}
  \caption{M4\_T3}
\end{subfigure}
\caption{M4 模態之 Q4 與 T3 對照。}
\end{figure}

\section{初步觀察}

由圖可見，部分候選模態具有明顯局部震盪與鋸齒狀變形，顯示僅以 $k$ 值接近 $k_{\mathrm{exact}}$ 作為模態篩選條件並不足夠。後續可加入模態形狀判斷，例如與 $\sin(m\pi x/a)\sin(n\pi y/b)$ 的相關性比較，以排除高頻或轉角主導的數值模態。

\end{document}
